\section{The Wound of Death}

Death separates the soul from the body. It tears apart the unity of the human being and ends the familiar presence through which we have known one another. The hand can no longer answer our hand. A voice whose smallest inflections were familiar has become silent. The body is still the body of the person we love, but we cannot restore the living exchange by which our love was ordinarily expressed. The Church honors that body and prays for that person because she knows both the reality of the wound and the promise of its healing.

At Lazarus's tomb the Lord meets two sisters grieving for their brother. Martha and Mary both tell Jesus that, had he been there, Lazarus would not have died. Martha professes her faith in the resurrection at the last day, and Jesus draws that faith toward his own person: he is the resurrection and the life. Then he goes to the tomb, weeps, and calls Lazarus forth (Jn 11:17--44). His tears show his love for Lazarus and his sorrow among the bereaved, even as he comes to restore their brother to them.

Our sorrow may therefore be brought to him without first being made eloquent or serene. An empty place at table, the absence of a companion, or a recollection that suddenly interrupts an ordinary day can become the occasion of prayer. We may be unable to explain our grief even to ourselves. Christ knows it before we can explain it. Martha and Mary approached him with their loss; we too may speak to him through tears and trust his promise of resurrection.

Death also belongs to the history of sin. Genesis presents the return to dust within the catastrophe of humanity's disobedience (Gen 2:17; 3:19). Paul contrasts Adam's trespass with Christ's obedience and names death the wages of sin (Rom 5:12--21; 6:23). The Catechism receives this teaching in its distinction between the natural mortality of the human body and the gift by which God had destined humanity not to die (CCC 1006--1008). Our bodily nature is corruptible; immunity from death belonged to the original divine gift. Death's familiarity within our present condition does not make it the fulfillment for which the Creator made us.

This teaching concerns the condition of the human family. It does not authorize us to measure an individual's guilt by the timing or painfulness of that person's death. When Christ is told of violent deaths and a fallen tower, he refuses to identify the victims as greater sinners than their neighbors. He directs his hearers toward their own repentance (Lk 13:1--5). His answer requires us to examine our sins. Those who suffer need our help, and their calamity gives us no knowledge of their guilt.

The death of a child, an unexpected death, or a death after long suffering can expose how little we control. We need not describe such a loss as though its good were immediately apparent. The promise of resurrection concerns something God will do. It does not demand that the bereaved discover, within every detail of their present pain, a reason that would make the pain welcome. We can trust the crucified and risen Lord while still asking why this person died and why those who loved him must suffer the separation.

The Second Vatican Council speaks of the human person's recoil before death and the seed of eternity within us (\emph{Gaudium et spes} 18). Christ conquers death and calls the whole person to communion with God. Those who survive us may remember us with love, but our life before God continues even when they can no longer remember. We may therefore pray with confidence for the dead whose names and histories have been lost to us: God still knows each one.

